\documentclass[11pt]{article}
\usepackage[margin=1in]{geometry}
\usepackage{amsmath,amssymb,amsthm,lmodern}
\usepackage[colorlinks=true,linkcolor=blue,urlcolor=blue]{hyperref}
\newtheorem{theorem}{Theorem}
\newtheorem{lemma}{Lemma}
\setlength{\emergencystretch}{2em}
\title{A formal proof of the graph-normalization adjacency inequality}
\author{Myeongjin Kang}
\date{25 October 2026}
\begin{document}
\maketitle
\begin{abstract}
We prove the graph-normalization adjacency inequality for every finite simple graph and every strictly positive weight vector. Connectivity is unnecessary. The proof is a weighted sum of the elementary inequality $2ab\leq a^2+b^2$ and applies to any symmetric nonnegative matrix whose action on the initial vector is strictly positive. Lean verifies the entire matrix statement. This is a formal verification of a standard majorization consequence; no priority claim is made.
\end{abstract}
\section{Statement}
Let $A$ be the adjacency matrix of a finite simple undirected graph, put $B=A+I$, and let $x_i>0$ for every vertex. Set
\[
 s_i=(Bx)_i,\qquad y_i=\frac{x_i}{s_i}.
\]
The diagonal entries of $B$ are one, so $s_i\geq x_i>0$.
\begin{theorem}\label{thm:graph}
For every such graph and vector,
\[
 y^T(A+I)y\leq \lVert y\rVert_1.
\]
\end{theorem}
This is Conjecture 8.1, equation (19), of Guigues \cite{original}, the target of MathDB 349070. The source assumes the graph connected; the theorem above removes that assumption. The conclusion is solely this inequality, not convergence of the normalization iteration.
\section{A symmetric-matrix inequality}
\begin{lemma}\label{lem:weighted}
Let $B$ be a finite symmetric matrix with nonnegative entries, and let $x_i\geq0$ and $a_i\in\mathbb R$. Writing $s_i=\sum_jB_{ij}x_j$, one has
\[
 \sum_{i,j}B_{ij}(x_i a_i)(x_j a_j)\leq\sum_i x_i a_i^2s_i.
\]
\end{lemma}
\begin{proof}
Multiply $2a_i a_j\leq a_i^2+a_j^2$ by $B_{ij}x_i x_j\geq0$ and sum over all ordered pairs. Symmetry and exchange of the finite summation indices give
\begin{align*}
2\sum_{i,j}B_{ij}x_i x_j a_i a_j
&\leq\sum_{i,j}B_{ij}x_i x_j a_i^2+\sum_{i,j}B_{ij}x_i x_j a_j^2\\
&=2\sum_i x_i a_i^2\sum_j B_{ij}x_j.
\end{align*}
Dividing by two proves the lemma.
\end{proof}
\begin{proof}[Proof of Theorem~\ref{thm:graph}]
More generally, suppose $B$ satisfies Lemma~\ref{lem:weighted} and all $s_i>0$. Substitute $a_i=1/s_i$. Then $x_i a_i=y_i$ and $x_i a_i^2s_i=y_i$, so the lemma gives $y^TBy\leq\sum_i y_i$. For a graph, $B=A+I$ satisfies all hypotheses. Since $y_i>0$, its sum is exactly $\lVert y\rVert_1$.
\end{proof}
\section{Formal correspondence and prior status}
The Lean source is \nolinkurl{formal/GraphNormalization.lean}. Its general finite-sum lemma is \nolinkurl{symmetric_energy_le}; the substitution is \nolinkurl{normalized_energy_le}. The explicit graph kernel equals mathlib's adjacency matrix plus the identity by \nolinkurl{closedKernel_eq_adj_add_one}. The theorem \nolinkurl{mathdb349070} states the result using mathlib matrix-vector multiplication and dot product, with the $\ell^1$ norm written as the sum of absolute values. Positivity of every denominator and the absolute-value equality are proved internally.

The diagonal majorization bound is standard and appears explicitly in Appendix B of Guigues's later paper \cite{later}. Its displayed global-dominance inequality, with $y^k=x$ and test vector $y_i=x_i/(Bx)_i$, already yields the conclusion above. Our contribution is the direct elementary argument and its complete formal certificate; no new-to-literature assertion is made. The later paper's additional convergence claims are neither used nor verified here.

The release contains the independent review, dependency pins, and \nolinkurl{verify_outcome.py}. Run the verifier on the extracted package with Lean 4.34.1 and the pinned mathlib dependencies. All eight local definitions and theorems are checked, with only \nolinkurl{propext}, \nolinkurl{Classical.choice}, and \nolinkurl{Quot.sound} permitted. The formal proof covers arbitrary finite graphs, including disconnected graphs; it does not use numerical testing or an assumed matrix inequality.
\begin{thebibliography}{9}
\bibitem{original} Laurent Guigues. \emph{Concerning Iterative Graph Normalization and Maximum Weight Independent Sets}. arXiv:2012.07764v1, 2020. Conjecture 8.1, equation (19). \url{https://arxiv.org/html/2012.07764v1}.
\bibitem{later} Laurent Guigues. \emph{Graph Normalization: Fast Binarizing Dynamics for Differentiable MWIS}. arXiv:2605.05330v1, 2026. Appendix B, global dominance. \url{https://arxiv.org/html/2605.05330v1}.
\end{thebibliography}
\end{document}
